\ifdefined\gkver\else\def\gkver{student}\fi
\documentclass[\gkver]{gaokaozhenti}
\begin{document}

\chapter{2009年海南}

\begin{enumerate}
\item
%% number: 1
%% paperName: 2009年海南高考第 1 题 3 分
%% typeId: 1
%% type: 单选题
%% chapter: 力物体的平衡
%% point: 力的合成
%% method: 无
%% score: 3
%% degree: 650
%% duplicateId: 0
%% body:
两个大小分别为 $F_{1}$ 和 $F_{2}$（$F_{2}<F_{1}$）的力作用在同一质点上，它们的合力的大小 $F$ 满足\xzanswer{C}

\fourchoices[answer=C]
{$F_{2}\leq F\leq F_{1}$}
{$\frac{F_{1}-F_{2}}{2}\leq F\leq\frac{F_{1}+F_{2}}{2}$}
{$F_{1}-F_{2}\leq F\leq F_{1}+F_{2}$}
{$F_{1}^{2}-F_{2}^{2}\leq F\leq F_{1}^{2}+F_{2}^{2}$}

%% answer: C
%% memo:
\memoanswer{
共点的两个力合成，同向时最大为 $F_{1}+F_{2}$，反向时最小为 $F_{1}-F_{2}$。
}

\item
%% number: 2
%% paperName: 2009年海南高考第 2 题 3 分
%% typeId: 1
%% type: 单选题
%% chapter: 磁场
%% point: 左手定则
%% method: 无
%% score: 3
%% degree: 650
%% duplicateId: 0
%% body:
一根容易形变的弹性导线，两端固定。导线中通有电流，方向如图中箭头所示。当没有磁场时，导线呈直线状态；当分别加上方向竖直向上、水平向右或垂直于纸面向外的匀强磁场时，描述导线状态的四个图示中正确的是\xzanswer{D}

\fourchoices[answer=D, ispicture=true, h=2.4cm]
{figs/02a.png}
{figs/02b.png}
{figs/02c.png}
{figs/02d.png}

%% answer: D
%% memo:
\memoanswer{
匀强磁场竖直向上、和导线平行，导线受到安培力为 0，A 错；匀强磁场水平向右，根据左手定则可知导线受到安培力向里，B 错；匀强磁场垂直纸面向外，由左手定则可知导线受到安培力水平向右，C 错、D 对。
}

\item
%% number: 3
%% paperName: 2009年海南高考第 3 题 3 分
%% typeId: 1
%% type: 单选题
%% chapter: 力物体的平衡
%% point: 物体系平衡
%% method: 整体与隔离
%% score: 3
%% degree: 450
%% duplicateId: 0
%% body:
两刚性球 a 和 b 的质量分别为 $m_{a}$ 和 $m_{b}$、直径分别为 $d_{a}$ 和 $d_{b}$（$d_{a}>d_{b}$）。将 a、b 球依次放入一竖直放置、内径为 $d$（$d_{a}<d<d_{a}+d_{b}$）的平底圆筒内，如图所示。设 a、b 两球静止时对圆筒侧面的压力大小分别为 $f_{1}$ 和 $f_{2}$，筒底所受的压力大小为 $F$。已知重力加速度大小为 $g$。若所以接触都是光滑的，则\xzanswer{A}

\onepicture[w=3.0cm, align=right]{figs/03.png}

\fourchoices[answer=A]
{$F=(m_{a}+m_{b})g$，$f_{1}=f_{2}$}
{$F=(m_{a}+m_{b})g$，$f_{1}\neq f_{2}$}
{$m_{a}g<F<(m_{a}+m_{b})g$，$f_{1}=f_{2}$}
{$m_{a}g<F<(m_{a}+m_{b})g$，$f_{1}\neq f_{2}$}

%% answer: A
%% memo:
\memoanswer{
对两刚性球 a 和 b 整体分析，竖直方向平衡可知 $F=(m_{a}+m_{b})g$、水平方向平衡有 $f_{1}=f_{2}$。
}

\item
%% number: 4
%% paperName: 2009年海南高考第 4 题 3 分
%% typeId: 1
%% type: 单选题
%% chapter: 电磁感应
%% point: 楞次定律推广
%% method: 无
%% score: 3
%% degree: 450
%% duplicateId: 0
%% body:
一长直铁芯上绕有一固定线圈 M，铁芯右端与一木质圆柱密接，木质圆柱上套有一闭合金属环 N，N 可在木质圆柱上无摩擦移动。M 连接在如图所示的电路中，其中 $R$ 为滑线变阻器，$E_{1}$ 和 $E_{2}$ 为直流电源，S 为单刀双掷开关。下列情况中，可观测到 N 向左运动的是\xzanswer{C}

\onepicture[w=3.6cm, align=right]{figs/04.png}

\fourchoices[answer=C]
{在 S 断开的情况下，S 向 a 闭合的瞬间}
{在 S 断开的情况下，S 向 b 闭合的瞬间}
{在 S 已向 a 闭合的情况下，将 R 的滑动头向 c 端移动时}
{在 S 已向 a 闭合的情况下，将 R 的滑动头向 d 端移动时}

%% answer: C
%% memo:
\memoanswer{
在 S 断开的情况下 S 向 a（b）闭合的瞬间 M 中电流瞬时增加、左端为磁极 N（S）极，穿过 N 的磁通量增加，根据楞次定律阻碍磁通量变化可知 N 环向右运动，A、B 均错；在 S 已向 a 闭合的情况下将 $R$ 的滑动头向 c 端移动时，电路中电流减小，M 产生的磁场减弱，穿过 N 的磁通量减小，根据楞次定律阻碍磁通量变化可知 N 环向左运动，D 错、C 对。
}

\item
%% number: 5
%% paperName: 2009年海南高考第 5 题 3 分
%% typeId: 1
%% type: 单选题
%% chapter: 电路
%% point: 电容
%% method: 无
%% score: 3
%% degree: 650
%% duplicateId: 0
%% body:
一平行板电容器两极板间距为 $d$、极板面积为 $S$，电容为 $\frac{\varepsilon_{0}S}{d}$，其中 $\varepsilon_{0}$ 是常量。对此电容器充电后断开电源。当增加两板间距时，电容器极板间\xzanswer{A}

\fourchoices[answer=A]
{电场强度不变，电势差变大}
{电场强度不变，电势差不变}
{电场强度减小，电势差不变}
{电场强度较小，电势差减小}

%% answer: A
%% memo:
\memoanswer{
平行板所带电荷量为 $Q$、两板间电压为 $U$，有 $C=\frac{Q}{U}$、$C=\frac{\varepsilon_{0}S}{d}$、两板间匀强电场的场强 $E=\frac{U}{d}$，可得 $E=\frac{Q}{\varepsilon_{0}S}$。电容器充电后断开，电容器电荷量 $Q$ 不变，则 $E$ 不变。根据 $C=\frac{\varepsilon_{0}S}{d}$ 可知 $d$ 增大、$C$ 减小，又根据 $C=\frac{Q}{U}$ 可知 $U$ 变大。
}

\item
%% number: 6
%% paperName: 2009年海南高考第 6 题 3 分
%% typeId: 1
%% type: 单选题
%% chapter: 曲线运动
%% point: 人造地球卫星
%% method: 无
%% score: 3
%% degree: 450
%% duplicateId: 0
%% body:
近地人造卫星 1 和 2 绕地球做匀速圆周运动的周期分别为 $T_{1}$ 和 $T_{2}$，设在卫星 1、卫星 2 各自所在的高度上的重力加速度大小分别为 $g_{1}$、$g_{2}$，则\xzanswer{B}

\fourchoices[answer=B]
{$\frac{g_{1}}{g_{2}}=\left(\frac{T_{1}}{T_{2}}\right)^{4/3}$}
{$\frac{g_{1}}{g_{2}}=\left(\frac{T_{2}}{T_{1}}\right)^{4/3}$}
{$\frac{g_{1}}{g_{2}}=\left(\frac{T_{1}}{T_{2}}\right)^{2}$}
{$\frac{g_{1}}{g_{2}}=\left(\frac{T_{2}}{T_{1}}\right)^{2}$}

%% answer: B
%% memo:
\memoanswer{
卫星绕天体做匀速圆周运动，由万有引力提供向心力有 $\frac{GMm}{R^{2}}=m\left(\frac{2\pi}{T}\right)^{2}R$，可得 $\frac{T^{2}}{R^{3}}=K$ 为常数；由重力等于万有引力 $\frac{GMm}{R^{2}}=mg$，联立解得 $g=\frac{GM}{\sqrt[3]{\frac{T^{4}}{K^{2}}}}$，则 $g$ 与 $T^{\frac{4}{3}}$ 成反比。
}

\item
%% number: 7
%% paperName: 2009年海南高考第 7 题 4 分
%% typeId: 2
%% type: 多选题
%% chapter: 功和能
%% point: 功
%% method: 无
%% score: 4
%% degree: 450
%% duplicateId: 0
%% body:
一物体在外力的作用下从静止开始做直线运动，合外力方向不变，大小随时间的变化如图所示。设该物体在 $t_{0}$ 和 $2t_{0}$ 时刻相对于出发点的位移分别是 $x_{1}$ 和 $x_{2}$，速度分别是 $v_{1}$ 和 $v_{2}$，合外力从开始至 $t_{0}$ 时刻做的功是 $W_{1}$，从 $t_{0}$ 至 $2t_{0}$ 时刻做的功是 $W_{2}$，则\xzanswer{AC}

\onepicture[w=5.0cm, align=right]{figs/07.png}

\fourchoices[answer=AC]
{$x_{2}=5x_{1}$，$v_{2}=3v_{1}$}
{$x_{1}=9x_{2}$，$v_{2}=5v_{1}$}
{$x_{2}=5x_{1}$，$W_{2}=8W_{1}$}
{$v_{2}=3v_{1}$，$W_{2}=9W_{1}$}

%% answer: AC
%% memo:
\memoanswer{
根据 $F-t$ 图像面积意义和动量定理有 $mv_{1}=F_{0}t_{0}$、$mv_{2}=F_{0}t_{0}+2F_{0}t_{0}$，则 $v_{2}=3v_{1}$；应用位移公式知 $x_{1}=\frac{v_{1}}{2}t_{0}$、$x_{2}=\frac{v_{1}+v_{2}}{2}t_{0}+\frac{v_{1}}{2}t_{0}$，则 $x_{2}=5x_{1}$，B 错、A 对；在第一个 $t_{0}$ 内对物体用动能定理有 $W_{1}=\frac{mv_{1}^{2}}{2}$、在第二个 $t_{0}$ 内对物体应用动能定理有 $W_{2}=\frac{mv_{2}^{2}}{2}-\frac{mv_{1}^{2}}{2}$，则 $W_{2}=8W_{1}$，D 错、C 对。
}

\item
%% number: 8
%% paperName: 2009年海南高考第 8 题 4 分
%% typeId: 2
%% type: 多选题
%% chapter: 直线运动
%% point: 相遇问题
%% method: 无
%% score: 4
%% degree: 450
%% duplicateId: 0
%% body:
甲乙两车在一平直道路上同向运动，其 $v-t$ 图像如图所示，图中 $\triangle OPQ$ 和 $\triangle OQT$ 的面积分别为 $s_{1}$ 和 $s_{2}$（$s_{2}>s_{1}$），初始时，甲车在乙车前方 $s_{0}$ 处。\xzanswer{ABC}

\onepicture[w=4.0cm, align=right]{figs/08.png}

\fourchoices[answer=ABC]
{若 $s_{0}=s_{1}+s_{2}$，两车不会相遇}
{若 $s_{0}<s_{1}$，两车相遇 2 次}
{若 $s_{0}=s_{1}$，两车相遇 1 次}
{若 $s_{0}=s_{2}$，两车相遇 1 次}

%% answer: ABC
%% memo:
\memoanswer{
由图可知甲的加速度 $a_{1}$ 比乙的 $a_{2}$ 大，在达到速度相等的时间 $T$ 内两车相对位移为 $s_{1}$。若 $s_{0}=s_{1}+s_{2}$，速度相等时甲比乙位移多 $s_{1}<s_{0}$，乙车还没有追上，此后甲车比乙车快，不可能追上，A 对；若 $s_{0}<s_{1}$，乙车追上甲车时乙车比甲车快，因为甲车加速度大，甲车会再追上乙车，之后乙车不能再追上甲车，B 对；若 $s_{0}=s_{1}$，恰好在速度相等时追上、之后不会再相遇，C 对；若 $s_{0}=s_{2}$（$s_{2}>s_{1}$），两车速度相等时还没有追上，并且甲车快、更追不上，D 错。
}

\item
%% number: 9
%% paperName: 2009年海南高考第 9 题 4 分
%% typeId: 2
%% type: 多选题
%% chapter: 交流电
%% point: 变压器
%% method: 无
%% score: 4
%% degree: 450
%% duplicateId: 0
%% body:
一台发电机最大输出功率为 $4000\UkW$，电压为 $4000\UV$，经变压器 $T_{1}$ 升压后向远方输电。输电线路总电阻 $R=1\UkO$。到目的地经变压器 $T_{2}$ 降压，负载为多个正常发光的灯泡（$220\UV$、$60\UW$）。若在输电线路上消耗的功率为发电机输出功率的 $10\%$，变压器 $T_{1}$ 和 $T_{2}$ 的耗损可忽略，发电机处于满负荷工作状态，则\xzanswer{ABD}

\fourchoices[answer=ABD]
{$T_{1}$ 原、副线圈电流分别为 $10^{3}\UA$ 和 $20\UA$}
{$T_{2}$ 原、副线圈电压分别为 $1.8\times10^{5}\UV$ 和 $220\UV$}
{$T_{1}$ 和 $T_{2}$ 的变压比分别为 $1:50$ 和 $40:1$}
{有 $6\times10^{4}$ 盏灯泡（$220\UV$、$60\UW$）正常发光}

%% answer: ABD
%% memo:
\memoanswer{
输电线路上消耗的功率为 $P=400\UkW=I^{2}R$ 可知输电线上电流为 $I_{2}=20\UA$，根据原线圈 $P_{1}=U_{1}I_{1}$，可知 $I_{1}=10^{3}\UA$，A 对，$T_{1}$ 的变压比为 $I_{2}:I_{1}=1:50$；根据 $P=U_{2}I_{2}$，可知 $U_{2}=2\times10^{5}\UV$，输电线上电压 $U_{\text{线}}=I_{2}R=20000\UV$，则副线圈输入电压为 $U_{3}=U_{2}-U_{\text{线}}=1.8\times10^{5}\UV$，又灯泡正常发光、副线圈电压为 $220\UV$，B 对，$T_{2}$ 的变压比为 $U_{3}:220$，C 错；根据 $U_{3}I_{2}=60n$，解得 $n=6\times10^{4}$，D 对。
}

\item
%% number: 10
%% paperName: 2009年海南高考第 10 题 4 分
%% typeId: 2
%% type: 多选题
%% chapter: 电场
%% point: 电场叠加
%% method: 无
%% score: 4
%% degree: 250
%% duplicateId: 0
%% body:
如图，两等量异号的点电荷相距为 $2a$。M 与两点电荷共线，N 位于两点电荷连线的中垂线上，两点电荷连线中点到 M 和 N 的距离都为 $L$，且 $L\gg a$。略去 $\left(\frac{a}{L}\right)^{n}$（$n\geq2$）项的贡献，则两点电荷的合电场在 M 和 N 点的强度\xzanswer{AC}

\onepicture[w=4.2cm, align=right]{\includesvg{figs/10}}

\fourchoices[answer=AC]
{大小之比为 2，方向相反}
{大小之比为 1，方向相反}
{大小均与 $a$ 成正比，方向相反}
{大小均与 $L$ 的平方成反比，方向相互垂直}

%% answer: AC
%% memo:
\memoanswer{
如右下图所示，合电场在 M 和 N 点的强度分别为 $E_{1}=\frac{Kq}{(L-a)^{2}}-\frac{Kq}{(L+a)^{2}}=\frac{4Kqa}{L^{3}}$、$E_{2}=2\frac{Kq}{L^{2}+a^{2}}\times\frac{a}{\sqrt{L^{2}+a^{2}}}=\frac{2Kqa}{L^{3}}$，$E_{1}:E_{2}=2$；又 N 点处场强方向由 $+q$ 指向 $-q$，在 M 点的场强表现为 $+q$ 的点电荷、由 $-q$ 指向 $+q$。

\onepicture[w=5.0cm]{figs/10_an.png}
}

\item
%% number: 11
%% paperName: 2009年海南高考第 11 题 4 分
%% typeId: 3
%% type: 填空题
%% chapter: 曲线运动
%% point: 万有引力定律
%% method: 无
%% score: 4
%% degree: 650
%% duplicateId: 0
%% body:
在下面括号内列举的科学家中，对发现和完善万有引力定律有贡献的是\tkanswer[6cm]{第谷、开普勒、牛顿、卡文迪许}。（安培、牛顿、焦耳、第谷、卡文迪许、麦克斯韦、开普勒、法拉第）

%% answer:
%% memo:
\memoanswer{
第谷搜集记录天文观测资料、开普勒发现开普勒三定律、牛顿发现万有引力定律、卡文迪许测定万有引力常数。
}

\item
%% number: 12
%% paperName: 2009年海南高考第 12 题 4 分
%% typeId: 3
%% type: 填空题
%% chapter: 交流电
%% point: 交流电
%% method: 无
%% score: 4
%% degree: 450
%% duplicateId: 0
%% body:
钳型表的工作原理如图所示。当通有交流电的导线从环形铁芯的中间穿过时，与绕在铁芯上的线圈相连的电表指针会发生偏转。由于通过环形铁芯的磁通量与导线中的电流成正比，所以通过偏转角度的大小可以测量导线中的电流。日常所用交流电的频率在中国和英国分别为 $50\UHz$ 和 $60\UHz$。现用一钳型电流表在中国测量某一电流，电表读数为 $10\UA$；若用同一电表在英国测量同样大小的电流，则读数将是\tkanswer[1.5cm]{12}$\UA$。若此表在中国的测量值是准确的，且量程为 $30\UA$；为使其在英国的测量值变为准确，应重新将其量程标定为\tkanswer[1.5cm]{25}$\UA$。

\onepicture[w=4.2cm, align=right]{figs/12.png}

%% answer:
%% memo:
\memoanswer{
根据 $\varphi\propto i$，$i=I_{m}\sin(2\pi ft)$，$\frac{\Delta\varphi}{\Delta t}\propto I_{m}2\pi f\cos(2\pi ft)$，说明线圈中的电动势有效值与频率成正比，根据欧姆定律可知电流与频率成正比，所以在英国用同一电表测量同样大小的电流的读数将是 $I=10\UA\times\frac{60}{50}=12\UA$；因为在英国电流值为标准的 $6/5$，需量程变为原来的 $5/6$，为 $25\UA$。
}

\item
%% number: 13
%% paperName: 2009年海南高考第 13 题 4 分
%% typeId: 4
%% type: 实验题
%% chapter: 直线运动
%% point: 游标卡尺和螺旋测微仪
%% method: 无
%% score: 4
%% degree: 650
%% duplicateId: 0
%% body:
某同学用游标卡尺和螺旋测微器分别测量一薄的金属圆片的直径和厚度，读出图中的示数，该金属圆片的直径的测量值为\tkanswer[2.4cm]{$1.240$}$\Ucm$，厚度的测量值为\tkanswer[2.4cm]{$1.682$（或 $1.683$）}$\Umm$。

\onepicture[w=8.0cm, align=right]{figs/13.png}

%% answer:
%% memo:
\memoanswer{
本题原卷及材料中未提供详解，待补充。
}

\item
%% number: 14
%% paperName: 2009年海南高考第 14 题 11 分
%% typeId: 4
%% type: 实验题
%% chapter: 电路
%% point: 测电动势与内阻
%% method: 无
%% score: 11
%% degree: 450
%% duplicateId: 0
%% body:
图 1 是利用两个电流表 $A_{1}$ 和 $A_{2}$ 测量干电池电动势 $E$ 和内阻 $r$ 的电路原理图。图中 S 为开关，$R$ 为滑动变阻器，固定电阻 $R_{1}$ 和 $A_{1}$ 内阻之和为 $10000\UO$（比 $r$ 和滑动变阻器的总电阻都大得多），$A_{2}$ 为理想电流表。

\onepicture[num=1, numstyle=arabic, showlabel=true, w=5.0cm, align=right]{figs/14a.png}

\begin{enumerate}
	\item
	按电路原理图在图 2 虚线框内各实物图之间画出连线。
	\item
	在闭合开关 S 前，将滑动变阻器的滑动端 c 移动至\tkanswer[2cm]{}（填“a 端”、“中央”或“b 端”）。
	\item
	闭合开关 S，移动滑动变阻器的滑动端 c 至某一位置，读出电流表 $A_{1}$ 和 $A_{2}$ 的示数 $I_{1}$ 和 $I_{2}$。多次改变滑动端 c 的位置，得到的数据为
	\begin{center}
	\begin{tabular}{|c|c|c|c|c|c|c|}
	\hline
	$I_{1}$（$\UmA$） & 0.120 & 0.125 & 0.130 & 0.135 & 0.140 & 0.145 \\
	\hline
	$I_{2}$（$\UmA$） & 480 & 400 & 320 & 232 & 140 & 68 \\
	\hline
	\end{tabular}
	\end{center}

	在图 3 所示的坐标纸上以 $I_{1}$ 为纵坐标、$I_{2}$ 为横坐标画出所对应的 $I_{1}-I_{2}$ 曲线。
	\item
	用所得曲线求的电源的电动势 $E=$\tkanswer[1.5cm]{}$\UV$，内阻 $r=$\tkanswer[1.5cm]{}$\UO$。（保留两位小数）
	\item
	该电路中电源输出的短路电流 $I_{m}=$\tkanswer[1.5cm]{}$\UA$。
\end{enumerate}
\twopicture[numA=2, numB=3, numstyle=arabic, w=5.0cm, align=right]
{figs/14b.png}
{figs/14c.png}

%% answer:
\jdanswer{
\begin{enumerate}
	\item
	（连线图略）
	\item
	b 端
	\item
	（作图略）
	\item
	$E\approx1.49\UV$，$r\approx0.60\UO$
	\item
	$I_{m}\approx2.5\UA$
\end{enumerate}
}

%% memo:
\memoanswer{
（1）实验前滑动变阻器接入电路电阻值最大，故滑动端 c 应移至 b 端。
（2）由图线上读出两组数值，代入 $E=I_{1}R_{1}+(I_{1}+I_{2})r$ 构成方程组联立求解 $E$ 和 $r$。
（5）短路电流 $I_{m}=E/r$。
}

\item
%% number: 15
%% paperName: 2009年海南高考第 15 题 9 分
%% typeId: 6
%% type: 计算题
%% chapter: 运动和力
%% point: 牛顿定律的应用
%% method: 无
%% score: 9
%% degree: 450
%% duplicateId: 0
%% body:
一卡车拖挂一相同质量的车厢，在水平直道上以 $v_{0}=12\Ums$ 的速度匀速行驶，其所受阻力可视为与车重成正比，与速度无关。某时刻，车厢脱落，并以大小为 $a=2\Umsq$ 的加速度减速滑行。在车厢脱落 $t=3\Us$ 后，司机才发觉并紧急刹车，刹车时阻力为正常行驶时的 3 倍。假设刹车前牵引力不变，求卡车和车厢都停下后两者之间的距离。

%% answer:
\jdanswer{$36\Um$}

%% memo:
\memoanswer{
设卡车的质量为 $M$，车所受阻力与车重之比为 $\mu$；刹车前卡车牵引力的大小为 $F$，卡车刹车前后加速度的大小分别为 $a_{1}$ 和 $a_{2}$。重力加速度大小为 $g$。由牛顿第二定律有
$F-2\mu Mg=0$ ①

$F-\mu Mg=Ma_{1}$ ②

$\mu Mg=Ma$ ③

$3\mu Mg=Ma_{2}$ ④

设车厢脱落后，$t=3\Us$ 内卡车行驶的路程为 $s_{1}$，末速度为 $v_{1}$，根据运动学公式有
$s_{1}=v_{0}t+\frac{1}{2}a_{1}t^{2}$ ⑤

$v_{1}=v_{0}+a_{1}t$ ⑥

$v_{1}^{2}=2a_{2}s_{2}$ ⑦

式中，$s_{2}$ 是卡车在刹车后减速行驶的路程。设车厢脱落后滑行的路程为 $s$，有 $v_{0}^{2}=2as$ ⑧

卡车和车厢都停下来后相距 $\Delta s=s_{1}+s_{2}-s$ ⑨

由①至⑨式得 $\Delta s=-\frac{v_{0}^{2}}{3a}+\frac{4}{3}v_{0}t+\frac{2}{3}at^{2}$。

代入题给数据得 $\Delta s=36\Um$。
}

\item
%% number: 16
%% paperName: 2009年海南高考第 16 题 10 分
%% typeId: 6
%% type: 计算题
%% chapter: 磁场
%% point: 带电粒子在磁场中的运动
%% method: 无
%% score: 10
%% degree: 450
%% duplicateId: 0
%% body:
如图，ABCD 是边长为 $a$ 的正方形。质量为 $m$、电荷量为 $e$ 的电子以大小为 $v_{0}$ 的初速度沿纸面垂直于 BC 边射入正方形区域。在正方形内适当区域中有匀强磁场。电子从 BC 边上的任意点入射，都只能从 A 点射出磁场。不计重力，求：

\begin{enumerate}
	\item
	次匀强磁场区域中磁感应强度的方向和大小；
	\item
	此匀强磁场区域的最小面积。
\end{enumerate}
\onepicture[w=3.2cm, align=right]{figs/16.png}

%% answer:
\jdanswer{
\begin{enumerate}
	\item
	$B=\frac{mv_{0}}{ea}$
	\item
	$\frac{\pi-2}{2}a^{2}$
\end{enumerate}
}

%% memo:
\memoanswer{
（1）设匀强磁场的磁感应强度的大小为 $B$。令圆弧 AEC 是自 C 点垂直于 BC 入射的电子在磁场中的运行轨道。电子所受到的磁场的作用力 $f=ev_{0}B$ ①，应指向圆弧的圆心，因而磁场的方向应垂直于纸面向外。圆弧 $\widehat{AEC}$ 的圆心在 CB 边或其延长线上。依题意，圆心在 A、C 连线的中垂线上，故 B 点即为圆心，圆半径为 $a$，按照牛顿定律有 $f=m\frac{v_{0}^{2}}{a}$ ②。联立①②式得 $B=\frac{mv_{0}}{ea}$ ③。

（2）由（1）中决定的磁感应强度的方向和大小，可知自 C 点垂直于 BC 入射的电子在 A 点沿 DA 方向射出，且自 BC 边上其它点垂直于 BC 入射的电子的运动轨道只能在 BAEC 区域中。因而，圆弧 AEC 是所求的最小磁场区域的一个边界。

为了决定该磁场区域的另一边界，我们来考察射中 A 点的电子的速度方向与 BA 的延长线交角为 $\theta$（不妨设 $0\leq\theta<\frac{\pi}{2}$）的情形。该电子的运动轨迹 $qpA$ 如右图所示。图中，圆弧 AP 的圆心为 O，$pq$ 垂直于 BC 边，由（1）知，圆弧 $\widehat{AP}$ 的半径仍为 $a$。在 D 为原点、DC 为 $x$ 轴、AD 为 $y$ 轴的坐标系中，P 点的坐标 $(x,y)$ 为 $x=a\sin\theta$ ④、$y=-a\cos\theta$ ⑤。

这意味着，在范围 $0\leq\theta\leq\frac{\pi}{2}$ 内，P 点形成以 D 为圆心、$a$ 为半径的四分之一圆周 $\widehat{AFC}$，它是电子做直线运动和圆周运动的分界线，构成所求磁场区域的另一边界。

因此，所求的最小匀强磁场区域是分别以 B 和 D 为圆心、$a$ 为半径的两个四分之一圆周 AEC 和 AFC 所围成的，其面积为 $S=2\left(\frac{1}{4}\pi a^{2}-\frac{1}{2}a^{2}\right)=\frac{\pi-2}{2}a^{2}$。

\onepicture[w=5.0cm]{figs/16_an.png}
}

\item
%% number: 17
%% paperName: 2009年海南高考第 17 题 8 分
%% typeId: 6
%% type: 计算题
%% chapter: 气体性质
%% point: 气体多过程
%% method: 无
%% score: 8
%% degree: 650
%% duplicateId: 0
%% body:
【模块 3-3 试题】（12 分）

\begin{enumerate}
	\item
	（4 分）下列说法正确的是\tkanswer[2cm]{ADEF}。（填入正确选项前的字母，每选错一个扣 2 分，最低得分为 0 分）

	\sixchoices[answer=ADEF]
	{气体的内能是分子热运动的动能和分子间的势能之和}
	{气体的温度变化时，其分子平均动能和分子间势能也随之改变}
	{功可以全部转化为热，但热量不能全部转化为功}
	{热量能够自发地从高温物体传递到低温物体，但不能自发地从低温物体传递到高温物体}
	{一定量的气体，在体积不变时，分子每秒平均碰撞次数随着温度降低而减小}
	{一定量的气体，在压强不变时，分子每秒对器壁单位面积平均碰撞次数随着温度降低而增加}
	\item
	（8 分）一气象探测气球，在充有压强为 $1.00\Uatm$（即 $76.0\UcmHg$）、温度为 $27.0\UdC$ 的氦气时，体积为 $3.50\Umc$。在上升至海拔 $6.50\Ukm$ 高空的过程中，气球内氦气逐渐减小到此高度上的大气压 $36.0\UcmHg$，气球内部因启动一持续加热过程而维持其温度不变。此后停止加热，保持高度不变。已知在这一海拔高度气温为 $-48.0\UdC$。求：
	\begin{enumerate}
		\item
		氦气在停止加热前的体积；
		\item
		氦气在停止加热较长一段时间后的体积。
	\end{enumerate}
\end{enumerate}

%% answer:
\jdanswer{
\begin{enumerate}
	\item
	ADEF
	\item
	（1）$V_{2}=7.39\Umc$；（2）$V_{3}=5.54\Umc$。
\end{enumerate}
}

%% memo:
\memoanswer{
在气球上升至海拔 $6.50\Ukm$ 高空的过程中，气球内氦气经历一等温过程。根据玻意耳—马略特定律有 $p_{1}V_{1}=p_{2}V_{2}$ ①，式中，$p_{1}=76.0\UcmHg$、$V_{1}=3.50\Umc$、$p_{2}=36.0\UcmHg$，$V_{2}$ 是在此等温过程末氦气的体积。由①式得 $V_{2}=7.39\Umc$ ②。

在停止加热较长一段时间后，氦气的温度逐渐从 $T_{1}=300\UK$ 下降到与外界气体温度相同，即 $T_{2}=225\UK$。这是一等压过程，根据盖—吕萨克定律有 $\frac{V_{2}}{T_{1}}=\frac{V_{3}}{T_{2}}$ ③，式中，$V_{3}$ 是在此等压过程末氦气的体积。由③式得 $V_{3}=5.54\Umc$ ④。
}

\item
%% number: 18
%% paperName: 2009年海南高考第 18 题 7 分
%% typeId: 6
%% type: 计算题
%% chapter: 机械波
%% point: 波长频率波速
%% method: 无
%% score: 7
%% degree: 650
%% duplicateId: 0
%% body:
【模块 3-4 试题】（12 分）

\begin{enumerate}
	\item
	（5 分）如图，一透明半圆柱体折射率为 $n=2$，半径为 $R$、长为 $L$。一平行光束从半圆柱体的矩形表面垂直射入，从部分柱面有光线射出。求该部分柱面的面积 $S$。
	\onepicture[w=4.0cm, align=right]{figs/18a.png}
	\item
	（7 分）有一种示波器可以同时显示两列波形。对于这两列波，显示屏上横向每格代表的时间间隔相同。利用此示波器可以测量液体中的声速，实验装置的一部分如图 1 所示：管内盛满液体，音频信号发生器所产生的脉冲信号由置于液体内的发射器发出，被接收器所接收。图 2 为示波器的显示屏。屏上所显示的上、下两列波形分别为发射信号与接收信号。若已知发射的脉冲信号频率为 $f=2000\UHz$，发射器与接收器的距离为 $s=1.30\Um$，求管内液体中的声速。（已知所测声速应在 $1300\sim1600\Ums$ 之间，结果保留两位有效数字。）
	\twopicture[numA=1, numB=2, numstyle=arabic, wA=7.2cm, wB=4.2cm, align=right]
	{figs/18b.png}
	{figs/18c.png}
\end{enumerate}

%% answer:
\jdanswer{
\begin{enumerate}
	\item
	$S=\frac{\pi}{3}RL$
	\item
	$v=1.4\times10^{3}\Ums$
\end{enumerate}
}

%% memo:
\memoanswer{
（1）半圆柱体的横截面如图所示，$OO'$ 为半径。设从 A 点入射的光线在 B 点处恰好满足全反射条件，由折射定律有 $n\sin\theta=1$ ①，式中，$\theta$ 为全反射临界角。由几何关系得 $\angle OO'B=\theta$ ②、$S=2RL\cdot\angle O'OB$ ③。代入题给条件得 $S=\frac{\pi}{3}RL$ ④。

\onepicture[w=4.2cm]{figs/18_an.png}

（2）设脉冲信号的周期为 $T$，从显示的波形可以看出，图 2 中横向每一分度（即两条长竖线间的距离）所表示的时间间隔为 $\Delta t=\frac{T}{2}$ ⑤，其中 $T=\frac{1}{f}$ ⑥。对比图 2 中上、下两列波形，可知信号在液体中从发射器传播至接收器所用的时间为 $t=(\Delta t)(2n+1.6)$ ⑦，其中 $n=0,1,2,\ldots$。液体中的声速为 $v=\frac{s}{t}$ ⑧。联立⑤⑥⑦⑧式，代入已知条件并考虑到所测声速应在 $1300\sim1600\Ums$ 之间，得 $v=1.4\times10^{3}\Ums$。
}

\item
%% number: 19
%% paperName: 2009年海南高考第 19 题 7 分
%% typeId: 6
%% type: 计算题
%% chapter: 原子和原子核
%% point: 质能方程
%% method: 无
%% score: 7
%% degree: 650
%% duplicateId: 0
%% body:
【模块 3-5 试题】（12 分）

\begin{enumerate}
	\item
	（5 分）已知：功率为 $100\UW$ 灯泡消耗的电能的 $5\%$ 转化为所发出的可见光的能量，光速 $c=3.0\times10^{8}\Ums$，普朗克常量 $h=6.63\times10^{-34}\UJcdots$，假定所发出的可见光的波长都是 $560\Unm$，计算灯泡每秒内发出的光子数。
	\item
	（7 分）钚的放射性同位素 $\ce{^{239}_{94}Pu}$ 静止时衰变为铀核激发态 $\ce{^{235}_{92}U^{*}}$ 和 $\alpha$ 粒子，而铀核激发态 $\ce{^{235}_{92}U^{*}}$ 立即衰变为铀核 $\ce{^{235}_{92}U}$，并放出能量为 $0.097\UMeV$ 的 $\gamma$ 光子。已知：$\ce{^{239}_{94}Pu}$、$\ce{^{235}_{92}U}$ 和 $\alpha$ 粒子的质量分别为 $m_{\mathrm{Pu}}=239.0521\Uu$、$m_{\mathrm{U}}=235.0439\Uu$ 和 $m_{\alpha}=4.0026\Uu$，$1\Uu=931.5\UMeVcSq$。
	\begin{enumerate}
		\item
		写出衰变方程；
		\item
		已知衰变放出的光子的动量可忽略，求 $\alpha$ 粒子的动能。
	\end{enumerate}
\end{enumerate}

%% answer:
\jdanswer{
\begin{enumerate}
	\item
	$n=1.4\times10^{19}\Usn$
	\item
	（1）$\ce{^{239}_{94}Pu -> ^{235}_{92}U + \alpha + \gamma}$；（2）$E_{\alpha}=5.034\UMeV$。
\end{enumerate}
}

%% memo:
\memoanswer{
（1）一波长为 $\lambda$ 的光子能量为 $E_{\gamma}=\frac{hc}{\lambda}$ ①。设灯泡每秒内发出的光子数为 $n$，灯泡电功率为 $P$，则 $n=\frac{kP}{E}$ ②，式中，$k=5\%$ 是灯泡的发光效率。联立①②式得 $n=\frac{kP\lambda}{hc}$ ③。代入题给数据得 $n=1.4\times10^{19}\Usn$ ④。

（2）衰变方程为 $\ce{^{239}_{94}Pu -> ^{235}_{92}U^{*} + \alpha}$、$\ce{^{235}_{92}U^{*} -> ^{235}_{92}U + \gamma}$，或合起来有 $\ce{^{239}_{94}Pu -> ^{235}_{92}U + \alpha + \gamma}$。

上述衰变过程的质量亏损为 $\Delta m=m_{\mathrm{Pu}}-m_{\mathrm{U}}-m_{\alpha}$ ⑤，放出的能量为 $\Delta E=c^{2}\Delta m$ ⑥。这能量是铀核 $\ce{^{235}_{92}U}$ 的动能 $E_{U}$、$\alpha$ 粒子的动能 $E_{\alpha}$ 和 $\gamma$ 光子的能量 $E_{\gamma}$ 之和，$\Delta E=E_{U}+E_{\alpha}+E_{\gamma}$ ⑦。由⑤⑥⑦式得 $E_{U}+E_{\alpha}=(m_{\mathrm{Pu}}-m_{\mathrm{U}}-m_{\alpha})c^{2}-E_{\gamma}$ ⑧。

设衰变后的铀核和 $\alpha$ 粒子的速度分别为 $v_{U}$ 和 $v_{\alpha}$，则由动量守恒有 $m_{U}v_{U}=m_{\alpha}v_{\alpha}$ ⑨。又由动能的定义知 $E_{U}=\frac{1}{2}m_{U}v_{U}^{2}$、$E_{\alpha}=\frac{1}{2}m_{\alpha}v_{\alpha}^{2}$ ⑩。由⑧⑨⑩式得 $\frac{E_{U}}{E_{\alpha}}=\frac{m_{\alpha}}{m_{U}}$，$E_{\alpha}=\frac{m_{U}}{m_{U}+m_{\alpha}}\left[(m_{\mathrm{Pu}}-m_{\mathrm{U}}-m_{\alpha})c^{2}-E_{\gamma}\right]$。

代入题给数据得 $E_{\alpha}=5.034\UMeV$。
}

\end{enumerate}

\end{document}
